\section{Problem Setup and Sealed Protocol}
\label{sec:setup}

We study a conditional question: once an unchanged parent unlearner reaches a
predeclared level of direct forgetting, which retained behaviors are damaged,
and can information available before that trajectory guide a limited repair
budget? The protocol therefore separates what is fixed before target
unlearning from what is observed afterward.

\subsection{Retained-Behavior Universe}


Let $\theta_0$ denote the frozen model before unlearning, let $\mathcal D_f$
denote the deletion request, and let $\mathcal C$ be an enumerated universe of
retained prompt--answer behaviors. For $x\in\mathcal C$, $\ell_x(\theta)$ is
mean answer-token negative log-likelihood, with prompt tokens masked. The
candidate universe need not equal a retain stream used by the parent
unlearner.

We split candidates by semantic group,
\begin{equation}
  \mathcal C
  =
  \mathcal C_{\mathrm{disc}}
  \mathbin{\dot\cup}
  \mathcal C_{\mathrm{audit}} .
  \label{eq:candidate-split}
\end{equation}
The discovery fold may fit empirical rank maps and supply eligible repair
examples. The audit fold is excluded from repair, hyperparameter selection,
stopping, and checkpoint selection. Its profile is computed at $\theta_0$ and
sealed before the target trajectory.

\subsection{Parent Trajectory, Forgetting Gate, and Damage}

A parent unlearner $u$ starts from $\theta_0$ and produces saved checkpoints
$\theta_1,\ldots,\theta_T$ using its fixed objective, data, budget, and random
seed. It receives neither the two-signal profile nor the repair pool. An
external gate selects the first saved checkpoint satisfying the common direct
forgetting condition:
\begin{equation}
  t^\dagger
  =
  \min\!\left\{
    t\in\mathcal T_{\mathrm{save}}:
    \mathsf{ForgetOK}(\theta_t;\mathcal D_f)=1
  \right\}.
  \label{eq:forgetting-gate}
\end{equation}
The gate, evaluation set, and save schedule are fixed before the target run.
A trajectory with no such checkpoint is recorded as non-reaching and has no
forgetting-conditional prediction or protection outcome.

Realized damage to retained behavior $x$ at checkpoint $t$ is
\begin{equation}
  d_t(x)=\ell_x(\theta_t)-\ell_x(\theta_0).
  \label{eq:realized-damage}
\end{equation}
Positive values mean that the retained answer became less likely. The primary
prediction target is $d_{t^\dagger}(x)$ on
$\mathcal C_{\mathrm{audit}}$. It is never used to choose a profile weight,
repair pool, parent checkpoint, or returned repair checkpoint.

\subsection{One Sealed Profile, Two Separate Tests}

At $\theta_0$, the method estimates Loss Susceptibility and
computes Representation Proximity for every candidate. Target-request-disjoint development
trajectories fix the numerical operating point, the scalar mixture weight, and
the repair protocol. The target profile then produces two precommitted
objects: a ranking of $\mathcal C_{\mathrm{audit}}$ for prospective prediction
and a Top-$K_p$ eligible pool from $\mathcal C_{\mathrm{disc}}$ for repair.

\emph{Prediction} compares the sealed audit ranking with
$d_{t^\dagger}$ only after the forgetting gate is reached.
\emph{Protection} starts from the same $\theta_{t^\dagger}$ and applies a
common guarded repair operator to pools chosen by different selectors. Its
confirmatory outcomes are measured on the untouched audit fold. Improvement
on the selected discovery pool is diagnostic only: a selector helps only if
held-out retained behavior improves while the common forgetting and utility
contract remains satisfied.

Figure~\ref{fig:profile-two-uses} summarizes the two coordinates and the two
separately evaluated uses of the sealed score.
